\chapter{Stress Tests}
\section*{What Is This Test?} Cardiac stress testing evaluates the heart during exercise or medicine-induced stress.\section*{Alternative Names} Exercise stress test; treadmill test; stress echocardiogram; nuclear stress test.\section*{Principle} ECG, symptoms, blood pressure, imaging, or perfusion are observed as cardiac workload increases.\section*{Why This Test Is Ordered} It evaluates exertional chest discomfort, breathlessness, exercise capacity, rhythm, or known coronary disease when clinically appropriate.\section*{Primary vs Secondary Test} It is a secondary cardiac test selected after risk assessment; it is not suitable for every symptom or asymptomatic person.\section*{How This Test Is Performed} The patient exercises on a treadmill or bicycle, or receives a medicine that increases cardiac workload while monitored.\section*{What Preparation Is Needed} Wear comfortable clothes and shoes; follow instructions about food, caffeine, and medicines. Report symptoms during testing.\section*{Understanding Results} ECG changes, imaging defects, symptoms, blood pressure, and capacity are interpreted together; false positives and negatives occur.\section*{Follow-up Tests} Coronary CT angiography, echocardiography, catheterization, rhythm monitoring, or cardiology review may follow.\section*{References} \begin{enumerate}\item MedlinePlus. Stress Tests.\item American Heart Association. Exercise stress testing resources.\end{enumerate}\clearpage\section*{Understanding Results and Follow-up}
The laboratory report should identify the specimen, method, units, reference interval or decision limit, and whether the result is positive, negative, abnormal, or inconclusive. Reference intervals vary with age, sex, pregnancy, specimen type, assay, and laboratory; a result outside a range is not automatically a diagnosis, and a result within a range does not always exclude disease. Discuss unexpected results with the ordering clinician, who can decide whether repeat or confirmatory testing, treatment, or referral is appropriate. Do not change medicines or delay urgent care based on an isolated result.
\section*{Regulatory Status and Authorization Date}This chapter describes a test category, not a specific commercial product. FDA, CDSCO, EU IVDR, and MHRA approval or registration dates therefore cannot be assigned without the assay manufacturer, product name, instrument, and intended use. For laboratory-developed tests, an individual agency approval date may not exist. See the Preface for the interpretation of regulatory status.
\clearpage
